\subsection{Integral local models for the selected lattices}
\label{subsec:cut-wach-models}

At unramified primes we need Wach modules realizing the signature lattices and their labeled jumps.

Keep the lift, the twisted coefficient place and its assigned
tuple fixed. At an odd rational prime \(p\) unramified in the auxiliary field \(E\), let \(F\) be a completion of \(E\) above \(p\). Use a finite extension of
the coefficient field containing all embeddings of \(F\) and
the first \(p\)-th roots of unity. Write \(O\) for its integers,
\(\varpi\) for its uniformizer and \(k_O\) for its residue field.
Write \(a\) for an ordinary coefficient embedding of
\(O_F\otimes_{\mathbb Z_p}O=\prod_a O\), and set
\(\rho(a)=a\circ\operatorname{Frob}_F\), with
\(\operatorname{Frob}_F\) the arithmetic Frobenius on \(O_F\).
The terminal Hodge label of the ordinary target \(a\) is
\(i=\rho(a)\). Expressing the same Frobenius arrow in the two
sets of labels gives
\begin{equation}\label{eq:wach-terminal-label-arrow}
 \begin{array}{rclcl}
 \text{ordinary embeddings}&\rho(a)&\xrightarrow{\ \varphi\ }&a,\\
 \text{terminal labels}&i^-=\rho(i)&\xrightarrow{\ \varphi\ }&i .
 \end{array}
\end{equation}

By \eqref{eq:wach-terminal-label-arrow} the ordinary factor carrying label \(i\) is
\(\rho^{-1}(i)\). This is the convention of
Paragraph~\ref{par:wach-level-models}. The de Rham evaluation at
the first root uses \(\varphi^{-1}\), so we transport the
crystalline coefficients of the target by inverse crystalline
Frobenius to read the filtration at \(i\).
The factors over distinct completions of \(E\) are
independent. The reflected arrays at conjugate embeddings use
the same predecessor permutation, since conjugation commutes
with the unramified Frobenius on the embeddings of \(E\).

\begin{auxlemma}[Integral Wach factors]\label{lem:covariant-wach-factor-free}
Let \(F/\mathbb Q_p\) be unramified, with \(p\) odd, and take the
coefficient ring \(O\) splitting its embeddings and containing the
first \(p\)-th roots of unity. For a stable free \(O\)-lattice \(T\)
in a crystalline representation, put
\[
 N(T)=D(T)\cap N(T[1/p]),\qquad N_a=e_aN(T)
\]
in the rational \'etale module, where \(e_a\) is the ordinary
idempotent of \(O_F\otimes_{\mathbb Z_p}O\). Each \(N_a\) is
free over \(O[[\pi]]\) of rank \(\operatorname{rk}_O T\).
\end{auxlemma}
\begin{proof}
Here \(D\) is the cyclotomic \'etale equivalence for integral
and finite \(\mathbb Z_p\)-representations. Use its
exactness, tensor compatibility and compatibility with scalar
actions \cite[Theorem A.3.4.3]{fontaine-etale}. Berger's Wach
comparison applies because \(F\) is unramified. It gives an
integral free lattice before extending coefficients, a basis of
\(D(T)\) after inverting the cyclotomic variable and completing,
trivial cyclotomic action at the origin and a Frobenius matrix
and inverse with at most powers of the first cyclotomic factor
as denominators. The ring action of \(O\) is carried along by
functoriality. For any coefficient idempotent \(e\) of \(O_F\otimes O\), the corresponding summand satisfies
\[
 eN\ \text{is finite over}\ O[[\pi]].
\]
Here \(N\) denotes the full integral Wach module under consideration. Before the coefficient
action is taken into account the integral Wach lattice is free
over \(O_F[[\pi]]\), so its reduction modulo \(\pi\) is
\(\mathbb Z_p\)-torsion-free. The direct summand \(N_a/\pi N_a\)
has the same property.

Since a power of \(\varpi\) is associated
to \(p\) in \(O\), it is also torsion-free over \(O\), hence free.
Choose lifts \(b_1,\ldots,b_s\) of an \(O\)-basis of this reduction. Nakayama gives a surjection
\[
 O[[\pi]]^s\to N_a
\]
with these images. Let \(K\) be the kernel of the indicated surjection. Reducing
a relation modulo \(\pi\) puts each coordinate in
\(\pi O[[\pi]]\). As \(N_a\) is \(\pi\)-torsion-free, division
by \(\pi\) gives another relation, so \(K\subseteq\pi K\).
Iterating gives
\[
       K\subseteq\bigcap_{k\ge1}\pi^k O[[\pi]]^s=0 .
\]
Thus the surjection is an isomorphism. The \'etale comparison identifies
\(s\) with the dimension of the coefficient representation.
\cite[Proposition III.4.2 and Corollary III.4.5]{berger-wach}
\end{proof}

For a finite extension \(O'/O\), choose an \(O\)-basis of
\(O'\). The integral and rational comparison modules of
\(T\otimes_O O'\), and their intersection, are the finite
direct sums determined by that basis. This yields
\[
       N(T\otimes_O O')=N(T)\otimes_O O'
\]
with its natural scalar action. The same finite direct sums on
continuous cochains commute with taking cohomology, and with the
images defining the integral finite condition. If \(e'\) is
the relative ramification index, tensoring a composition factor
\(O/\varpi O\) gives \(O'/\varpi O'\) of length \(e'\).
Dividing the respective lengths by \(\operatorname{length}_O(O/pO)\)
and \(\operatorname{length}_{O'}(O'/pO')\) therefore gives the
same normalized length. Faithful flatness
also detects the required inclusions over \(O\). We use
these comparisons only at the coefficient place already fixed.

\paragraph{The norm coordinate and the labeled heights.}
On each factor of Lemma~\ref{lem:covariant-wach-factor-free},
use the ordinary cyclotomic variable \(\pi\), with
\(\gamma(1+\pi)=(1+\pi)^{\chi(\gamma)}\) and
\(Q=\big((1+\pi)^p-1\big)/\pi\). With \(\Delta\) the tame
group of order \(p-1\), put
\[
     Z=\prod_{\delta\in\Delta}\delta(\pi).
\]
Since \(\delta(\pi)\equiv\chi(\delta)\pi\pmod{\pi^2}\),
we have \(Z=v\pi^{p-1}\) for a unit \(v\in O[[\pi]]^\times\).
For an invariant series \(f\) with first term \(a_n\pi^n\ne0\),
comparison on the \(n\)-th graded piece gives
\[
 a_n\chi(\delta)^n=a_n\quad(\delta\in\Delta),
 \qquad n=(p-1)m\quad(m\ge0).
\]
Subtract \(a_n v(0)^{-m}Z^m\) to raise the order. Iteration
in the complete ring, together with the invariance of \(Z\),
gives
\[
 O[[\pi]]^\Delta=O[[Z]].
\]
Lifting a basis of \(N(T)/\pi N(T)\) and averaging
it over \(\Delta\) gives an invariant basis on each factor.

In
this basis the tame-invariant lattice \(N_i\) is free over
\(O[[Z]]\) and gives the full ordinary Wach lattice after
extension to \(O[[\pi]]\). At \(\pi=\zeta^a-1\), with
\(\zeta\) primitive of order \(p\) and \(1\le a<p\), we have
\(Z=p\), by the product of the nonzero cyclotomic differences.
Weierstrass preparation thus makes \(Z-p\) associated to \(Q\)
over \(O[[\pi]]\). The constant term and the initial unit
coefficient give the norm parameter used in the lift comparison.

Choose a cyclotomic generator \(\gamma\) with
\[
 u=\chi(\gamma)=1+pc,\qquad c\in\mathbb Z_p^\times,\qquad v_p(u-1)=1.
\]
The quotient
\(\gamma(Z)/Z-1\) is a series in \(Z\) which vanishes at
\(Z=p\), since \(\gamma\) fixes each primitive root. Hence
\(\gamma(Z)=Z(1+(Z-p)A(Z))\) for an integral series \(A\).
The constant term of \(\gamma(Z)/Z\) is \(u^{p-1}\), from
the leading ordinary term of \(Z\). Consequently
\[
 1-pA(0)=u^{p-1}\equiv 1+(p-1)pc\pmod{p^2},
 \qquad A(0)\equiv c\pmod p .
\]
In the residue ring Frobenius substitutes \(Z^p\). In the
ordinary lifting calculation put \(X=-Z\), \(q=X+p\). With bars
denoting reduction in the residue ring the quadratic terms are
\begin{equation}\label{eq:wach-norm-cyclotomic-sign}
 \overline{\gamma(Z)}\equiv Z+\bar c Z^2\pmod{Z^3},\qquad
 \overline{\gamma(X)}\equiv X-\bar c X^2\pmod{X^3}.
\end{equation}
Taking \(u=1+p\) in \eqref{eq:wach-norm-cyclotomic-sign} gives
quadratic coefficient \(-1\) in \(X\). This change of variable also changes the sign
of the divided matrix \((\gamma-1)/Z\) on the constants.

In the terminal labels write \(P_i\) for the Frobenius matrix
from \(i^-\) to \(i\) and \(1+Z D_i\) for the cyclotomic matrix
on \(i\). If the Hodge jumps at \(i\) are \(d_{i,1},\ldots,d_{i,r}\),
the elementary divisors of \(P_i\) at a primitive root have
orders \(d_{i,a}\). For an unramified quadratic host and a
line of jumps \((0,1)\), the orders in ordinary target order
are \((1,0)\). The terminal indexing gives \((0,1)\).
\begin{lemma}[First-level heights in the terminal index]
\label{lem:first-level-terminal-heights}
At each terminal label \(i\) the elementary-divisor orders of the
first-level matrix \(P_i\) are the signed jumps \(d_{i,a}\). For
\(-M_i\le d_{i,a}\le H_i\) its two integral height bounds are
obtained by clearing \(P_i\) by \((Z-p)^{M_i}\) and \(P_i^{-1}\)
by \((Z-p)^{H_i}\). For nonnegative jumps each \(a\)-minor is
divisible by the power given by the sum of the \(a\) least
jumps. In particular the characteristic polynomial, including multiplicities, is
\[
 \det(T-\bar D_i(0))=\prod_a(T+\bar c\,d_{i,a}).
\]
\end{lemma}

\begin{proof}
For the general comparison, first take nonnegative jumps and
compare the Wach lattice to the constant crystalline lattice
over the ring of convergent series. Let \(C_i\) express the
Wach basis in the constant basis. At the root, in the
filtration-adapted basis, the Wach lattice has generators
\(s^{-d_{i,a}}\) times the constant generators. Thus \(C_i\)
has elementary-divisor orders \(-d_{i,a}\). The matrix \(U_i\)
of the Frobenius image of the source Wach basis in the target
constant basis is invertible over the root power series ring.
Indeed the root maps to \(\pi=0\), where the comparison has
orders zero, and crystalline Frobenius is invertible.
Write
\(C_i=A_i\operatorname{diag}(s^{-d_{i,a}})_a B_i\) with both
outer factors invertible in the root power series ring.

Since
\(C_iP_i=U_i\) and \(U_i\) is invertible there,
\[
       P_i=B_i^{-1}\operatorname{diag}(s^{d_{i,a}})_a
                               A_i^{-1}U_i .
\]
So its elementary-divisor orders are exactly \(d_{i,a}\).
The inverse Frobenius transport of the constant basis puts these
orders and the cyclotomic matrix at the same terminal label.
A common twist adds its rank-one order to all \(d_{i,a}\).
Cancelling the rank-one factor gives the same formula for signed jumps.
In the same root chart, with local variable
\(s=\log((1+\pi)/\zeta^a)\), the action is \(s\mapsto us\).
The constant crystalline basis is fixed, and the matrix on the
Wach lattice at the root has eigenvalues \(u^{-d_{i,a}}\)
with multiplicity. These are the first-level cases of
Berger's Propositions II.3.1, III.2.1 and III.2.2.
\cite{berger-wach,dousmanis-wach}\label{cor:cut-wach-models-1}

In particular, if \(-M_i\le d_{i,a}\le H_i\), then
\[
 (Z-p)^{M_i}P_i,\qquad (Z-p)^{H_i}P_i^{-1}
                \quad\text{have entries in }O[[Z]].
\]
For descent of these entrywise conditions, start with the
coarse Wach common denominator, a power of \(Q\). In ordinary
coordinates multiply by it and remove the excess factors by
monic division at every primitive root. A monic linear factor
with a root of positive valuation in \(O\) gives integral
division with remainder the value at that root. Its prime ideal
has quotient \(O\), so distinct factors can be removed with their
specified multiplicities. The factors have product \(Q\).
The resulting bounded entries are invariant, so they descend to
\(O[[Z]]\).

For nonnegative jumps, apply the same monic division to any \(a\)-minor \(D_a\) of \(P_i\). Ordering the jumps increasingly gives
\[
 (Z-p)^{\sum_{b=1}^{a}d_{i,(b)}}\mid D_a\quad\text{in }O[[Z]].
\]
At \(Z=p\) subtract the identity
from the cyclotomic matrix and divide by \(p\) before reduction.
Since \(Z=p\) reduces to \(0\), this gives
\[
     \det(T-\bar D_i(0))=\prod_a(T+\bar c d_{i,a})
\]
in the residue field. This uses that the coefficients
of the characteristic polynomial commute with reduction and
that \((u^{-d}-1)/p\equiv-cd\pmod p\). On a scalar line
the Frobenius ideal is \((Z-p)^d\), by the two height bounds
with the sole jump \(d\).
\end{proof}

\begin{lemma}[Height bounds on the induced residue lattices]
\label{lem:localized-induced-wach-heights}
Put \(R_0=k_O[[Z]]\) and \(K_0=k_O((Z))\). Let
\(L_i=N_i/\varpi N_i\) with the terminal labels of Lemma~\ref{lem:first-level-terminal-heights}. A Galois-stable flag in
\(T/\varpi T\) gives a flag in the corresponding \(K_0\)-factors,
of the same ranks. At each factor, intersecting with \(L_i\) and
projecting to a flag quotient gives a free \(R_0\)-lattice. Let \(B_i\) denote the Frobenius span from \(i^-\) to \(i\).
If \(Z^{H_i}L_i\subseteq B_i\subseteq Z^{-M_i}L_i\), write \(L_i'\) for any induced flag-subquotient lattice and \(B_i'\) for its Frobenius span. Then
\[
 Z^{H_i}L_i'\subseteq B_i'\subseteq Z^{-M_i}L_i'.
\] Each marked line is identified by its specified
crystalline rank-one comparison in the residual \'etale realization.
\end{lemma}

\begin{proof}
Exactness of the \'etale equivalence at multiplication by \(\varpi\)
identifies the reduced module with the module of \(T/\varpi T\).
The invariant field of the tame action on \(k_O((\pi))\) is \(K_0\).
Indeed \(Z\) has order \(p-1\), and division by this series gives
extension degree \(p-1\). The \(p-1\) substitutions are distinct by
their linear coefficients. Field descent preserves dimension, and
taking the tame invariants is exact by averaging. Also the map
from the Wach lattice to the \'etale module is \(\varpi\)-saturated:
for \(b\in D(T)\) one has
\[
 \varpi b\in D(T)\cap N(T[1/p])\quad\Longrightarrow\quad b\in D(T)\cap N(T[1/p]),
\]
since the rational Wach module is closed under \(\varpi^{-1}\). Reduction therefore gives the
specified inclusion of \(L_i\) into the residual \'etale realization.

For the maps of a flag, take inverse images of its subspaces in
\(T\). Each inverse image contains \(\varpi T\) and has the same
crystalline rationalization as \(T\). Functoriality and exactness
of the \'etale equivalence identify the images of these maps after
reduction. The given rank-one comparison determines the individual
marked line in this flag. In particular this determines its
position before forming any saturated power-series lattices.
\label{cor:cut-wach-models-2}

For a stable rational subspace \(W_i\) of a flag, set
\(S_i=L_i\cap W_i\). Its quotient \(L_i/S_i\) is torsion-free,
so both are free over \(R_0\). Denote by \(B_i\) the Frobenius
span of \(L_{i^-}\) in the terminal factor. The height conditions
reduce to
\[
    Z^{H_i} L_i\subseteq B_i\subseteq Z^{-M_i}L_i.
\]
Flat Frobenius base change and a basis adapted to \(S_{i^-}\)
identify the Frobenius span of \(S_{i^-}\) with \(B_i\cap W_i\).
Intersection of the displayed inclusions with \(W_i\) gives
\[
    Z^{H_i} S_i\subseteq B_i\cap W_i\subseteq Z^{-M_i}S_i.
\]
Taking images in a stable rational quotient preserves the two
inclusions. This gives the assertion for a subquotient by first
intersecting and then projecting.
\end{proof}

The ordinary row and column comparisons use actual surjections
from the matrix span at the retained point. Their rational sources
are subrepresentations of finite sums of the normalized family and
its dual twists. Strictness of crystalline comparison puts the
Hodge jumps of source, target and kernel in the same labeled
boxes. If a positive twist is used, the same shift applies to
source and target. For a finite class let \(T'\) be its extension of the trivial lattice by the target \(T_2\). Strictness gives the multiset identity
\[
 \mathrm{Jumps}_\tau(T')=\mathrm{Jumps}_\tau(T_2)\sqcup\{0\}.
\]
Thus the same boxes apply whenever they contain zero. These are the data for the integral Herr pairs.

At the complementary primes use the heights of
Lemma~\ref{lem:first-level-terminal-heights}. The local comparison
uses the full
crystalline lattices at the same marked point and the induced
models of Lemma~\ref{lem:localized-induced-wach-heights}. Its
reduced matrix on the full lattice is the divided cyclotomic
matrix computed above. The common containing lattice permits
the various integral presentations of a class. For the least
exponent comparison choose the full lattice of the characteristic-zero
eigensystem lifting the surviving weight. Its marked lines are
those of the actual residual congruence. The Wach bases may differ between comparisons. Keeping the terminal labels fixed identifies the residual extension being tested, and normalization returns its length bound to the original coefficient place.
